\section{הגדרת הבעיה}

\subsection{ניסוח הבעיה}
בגידולים אינטנסיביים על מצע מנותק, בית השורשים (Rootzone) הוא מערכת דינמית שקשה למדוד באופן רציף ומייצג. ערכי המוליכות החשמלית (EC) והחומציות (pH) מושפעים בו-זמנית מתנאי האקלים, מהמיקרו-אקלים בתוך החממה, ממשטר ההשקיה והדישון, ממצב הצמח ומהיסטוריית האירועים במצע. כתוצאה מכך, שינוי קטן באחד מגורמים אלו עשוי להוביל לשינוי משמעותי בסביבת השורשים.

הבעיה המרכזית במחקר נובעת מהפער בין הצורך בבקרה רציפה ומדויקת של בית השורשים לבין מגבלות אמצעי הניטור הקיימים כיום:

\begin{itemize}
    \item \textbf{מגבלת הניטור המקוטע:} בדיקות ידניות או בדיקות מעבדה של נקז/מצע מספקות מדידה מדויקת יחסית, אך הן דלילות בזמן ומייצגות נקודת זמן מסוימת בלבד. בין דגימה לדגימה עלולים להתרחש שינויים ברמת המליחות או החומציות, שאינם מתועדים בזמן אמת.
    \item \textbf{מגבלת החיישנים הרציפים:} חיישנים המוצבים במצע יכולים לספק סדרת זמן רציפה, אך הם רגישים לבעיות כיול, סחיפה, רעש וייצוגיות מרחבית. מאחר שהחיישן מודד נקודה אחת בלבד, ייתכן שהקריאה אינה מייצגת את מצב כלל בית השורשים.
    \item \textbf{ניהול תגובתי:} בהיעדר יכולת חיזוי, המגדל נדרש לעיתים להגיב לחריגות רק לאחר שהן כבר התרחשו. מצב זה מקשה על שמירה יציבה של תנאי בית השורשים ועלול להוביל לשימוש לא מיטבי במים, בדשן ובתיקוני חומציות.
\end{itemize}

לכן, הבעיה המחקרית היא כיצד ניתן להעריך את מצב בית השורשים בפרקי הזמן שבין המדידות הישירות, ולחזות את ערכי ה-pH וה-EC העתידיים על בסיס מידע זמין ממקורות אחרים.

\subsection{שאלת המחקר}
שאלת המחקר המרכזית היא:

\begin{quote}
\itshape
האם ניתן לפתח מודל למידת מכונה המסוגל לחזות את ערכי ה-pH וה-EC בבית השורשים בנקודת זמן עתידית, על בסיס מדידה נוכחית ידועה, נתוני אקלים ומיקרו-אקלים, נתוני השקיה ודישון, מצב הצמח והיסטוריית המדידות?
\end{quote}

מתוך שאלה זו נגזרות שאלות משנה:

\begin{itemize}
    \item \textbf{דיוק החיזוי:} מהי רמת הדיוק שניתן להשיג בחיזוי ערכי pH ו-EC, כאשר ההערכה מתבססת על מדדי השגיאה MAE ו-RMSE ועל מדד ההסבר \Rtwo{}, והאם דיוק זה מתאים לשימוש ככלי תומך החלטה?
    \item \textbf{תרומת המשתנים:} אילו קבוצות משתנים תורמות באופן המשמעותי ביותר לחיזוי השינויים בבית השורשים: תנאי אקלים, השקיה, דישון, מצב הצמח, מדידות עבר או משך הזמן בין המדידה הנוכחית לנקודת החיזוי?
    \item \textbf{יציבות לאורך זמן:} האם המודל שומר על ביצועים יציבים כאשר הוא נבחן בשיטת ולידציה קדימה בזמן, המדמה תרחיש שבו המודל מאומן על נתוני עבר ונדרש לחזות נקודות עתידיות?
    \item \textbf{יכולת הכללה:} האם המודל מצליח לשמור על רמת דיוק טובה גם בבדיקת הכללה על מרווחים מדולגים, שאינם משמשים לאימון המודל?
\end{itemize}

\subsection{השערת המחקר}
השערת המחקר היא שקיים קשר פיזיקלי-כימי ואגרונומי בין תנאי הסביבה, פעולות ההשקיה והדישון, מצב הצמח והשינויים בערכי ה-pH וה-EC בבית השורשים. קשר זה אינו בהכרח ליניארי, והוא עשוי לכלול השהיות, תגובות מצטברות ואירועים נקודתיים בעלי השפעה גבוהה.

בהתאם לכך, ההשערה היא שמודל למידת מכונה יוכל ללמוד את דפוסי התגובה של בית השורשים מתוך נתוני העבר, ולחזות את ערכי ה-pH וה-EC בנקודות זמן עתידיות. בפרט, משוער כי מודל היברידי, המשלב בין \GB{} ללמידת דפוסים מורכבים לבין רכיב ליניארי להתמודדות עם מגמות ואירועים חריגים, יוכל לשמש כ־\SoftSensor{} אמין יחסית, המשלים את המדידות הידניות ומצמצם את פערי המידע ביניהן.
